\section{結論}
\label{sec:conclusion}

本稿は、企業が標準の非互換性を私的に迂回できる能力が、正式な標準化coalitionのstabilityにどのように影響するかを問う。答えは、正式なcoalition membershipを変えない場合でも、私的互換性が政治的に重要となり得るということである。企業は政府が正式分割を選んだ後で追加標準をサポートするかを選び、政府はその私的応答が生み出す継続結果を評価する。したがって、私的互換性投資は正式な制度境界に付随する経済価値を変え得る。本モデルでは、このfeedbackが、企業の内生的な互換性選択とそれに続くCournot competitionという標準的なindustrial-organization marginを、地域標準化と国際標準化に対する政府の誘因へ結び付ける。

中心的メカニズムは選択的侵食である。私的迂回がない場合、地域標準化連合の加盟国はcoalition加盟国市場で純排除余剰を得るが、域外国市場では相互的不利益を負う。域外国企業だけがブロック標準を採用すると、加盟国市場の構成要素が消える。域外国企業はcoalition市場でもはや排除されないため、そこでの加盟国の消費者余剰と国内企業利益のブロックは完全互換のcounterpartへ移る。しかし、加盟国は域外国標準を採用していないため、域外国市場でのcoalition加盟国の不利益は残る。したがって、私的適応は国際標準化を再現することも、正式分割を直接変更することもなく、地域的排除を支えた利得構成要素を選択的に取り除く。モデルが特定するパラメータ領域では、私的互換性はこのため、正式調和のmarket-access roleの一部を技術的に代替しながら、政府の国際標準化への相対的誘因を高めるという限定された意味で政治的補完として作用し得る。

Theorem~\ref{thm:formal-coalition-stability} は、この継続利得メカニズムを正式coalition stabilityへ変換する。すべての $(c,v)\in\Omega_0$ について、compatibility fixed costが高く $F>2T_A$ なら、地域標準または各国別標準の下で私的採用は生じず、strict-blocking stable setは3つすべての対称な地域標準化連合を含む。これに対し、固定費用が中間領域 $F_L<F<2T_A$ にあると、各地域連合の域外国企業はブロック標準を採用する一方、加盟企業は相互採用しない。この域外国企業のみの迂回によって地域的取り決めを支えていた加盟国固有の優位が消え、同時に域外国も国際標準化を厳格に選好する。grand coalitionはその結果すべての地域標準化連合を厳格にblockし、本稿が維持するcoalition-stability conceptの下で国際標準化が一意にstableとなる。この逆転はproduct-market parameterのopen setと、非退化なcompatibility cost interval上で起こり、threshold-equality resultではない。また、本稿はすべての $F$ にわたるmonotonicityを推論しない。より低コストで生じるmultistandarding patternはこのheadline comparisonの外にある。

副次的な市場規模結果は、Main mechanismを生み出すことなくpartner identityを明確化する。Proposition~\ref{prop:market-size-partner-selection} は、私的採用がないことを条件として、政府が、他の点では対称な2つのpotential regional partnerのうち大きい方を選好することを示す。市場質量をoutsider treatmentからcoalition-member treatmentへ移す価値が $A-C=T_U>0$ だからである。$m_1=m_2=1$、$m_3=1-\delta$ の下では、大きい2国は互いを小さい国より選好し、小さい国は両者の間で無差別である。これはpartner-ranking resultであってcoalition-stability theoremではなく、Theorem~\ref{thm:formal-coalition-stability} は市場規模の非対称性を必要としない。より広く見ると、本モデルのindustrial-organization contributionはfirm conductからinstitutional incentiveへのfeedbackである。標準が生産物市場競争を形作り、私的互換性の選択が利益と国民厚生を変え、その継続利得がblocking incentiveを決める。したがって、standards economics上の含意は、互換性がmarket accessとcompetitionだけでなく、正式なstandard-setting arrangementの政治価値にも影響するということである。coalition-formationの観点からは、post-formation private responseが内生的なら、formal membershipだけではcoalition payoffは固定されない。

制度的含意は、政策処方より意図的に狭い。互換性が非互換標準の技術的帰結を小さくする場合でも、私的互換性とpublic harmonizationが政治的代替になるとは限らない。design、multi-standard production、certification、その他のcompatibility investmentの費用変化はformal arrangement間のrentを再配分し、地域的調和とより広い調和に対する政府支持を変え得る。ただし本稿が特徴付けるのはcoalition stabilityであり、global social optimalityではない。中間領域で国際標準化が一意にstableとなるという結果を、world welfareを最大化するという結果として読むべきではなく、instabilityを地域的取り決めのinefficiencyと読むべきでもない。このメカニズムはむしろ今後の実証研究に問いを示す。すなわち、private interoperability costの低下が、exclusive regional arrangementとbroader standards arrangementに対する政治的支持の変化と関連するかどうかである。

正式結果は、1国1企業、binary compatibility adoption、外生的な固定採用費用、transferのないexclusive-membership partition gameにおけるstrict blockingを備えた、意図的にparsimoniousな3か国Cournot environmentで導出される。より大きなcoalition network、bargainingとside payment、alternative oligopoly competition、endogenous compatibility investmentは、採用誘因またはcoalition deviationを変える可能性があり、別の分析を必要とする。自然な次の2つの方向は、より豊かなcoalition networkで選択的侵食が維持されるかを検討することと、技術適応とformal membershipの区別を維持したまま企業のcompatibility investmentを内生化することである。現在のモデル内での中心的な教訓はより明確である。私的互換性は、企業が標準をまたいでどのように競争するかを変えるだけでなく、政府がどの正式な標準化の取り決めを維持する意思を持つかを決める継続利得を変えるために重要である。
